\documentclass{article}

\usepackage{amsmath}
\usepackage{amssymb}
\usepackage{graphicx}
\usepackage{booktabs}
\usepackage{hyperref}
\usepackage{natbib}
\usepackage{xcolor}

\newcommand{\TODO}[1]{\textcolor{red}{TODO: #1}}

\title{Persistent Memory Turns Prompt Injection into Durable Agent Compromise}
\author{Debeshee Das}
\date{}

\begin{document}

\maketitle

\begin{abstract}
Large language model (LLM) agents augmented with tools and persistent memory enable personalized, cross-session assistance. However, this architectural shift introduces a new class of stateful vulnerabilities: a single malicious tool output (e.g., an email) can implant a latent payload in an agent’s memory that later triggers covert data exfiltration under benign user behavior.

We formalize and empirically evaluate \emph{persistent data-exfiltration attacks} in realistic single-user deployments. Our attacks require only a one-time indirect injection, survive across long sequences of benign sessions (up to 100 evaluated sessions), and generalize to held-out trigger instances within a domain. We instantiate an email assistant across four memory backends—sliding-window context, retrieval-augmented generation (RAG), explicit human-editable memory, and a state-of-the-art agentic memory system (Mem0)—and evaluate four defenses: user-prompt-only indexing, disabling indexing after untrusted tool outputs, length-limited memory, and a provable information-flow policy with a non-interference guarantee.

Using an OpenEvolve-based automated red-teaming pipeline, we show that heuristic mitigations reduce but do not eliminate risk, while the information-flow policy prevents exfiltration empirically at measurable utility cost. We further introduce a capability-aware utility evaluation framework that decomposes agent tasks into clearly defined capability classes, enabling transparent security–utility tradeoff analysis. Our results demonstrate that persistent memory transforms prompt injection from a transient model failure into a durable architectural compromise unless principled information-flow boundaries are enforced.
\end{abstract}

\section{Introduction}

LLM agents increasingly combine language models with tools (e.g., email, web browsing, file access) and persistent memory to provide long-term personalized assistance. While these capabilities significantly enhance utility, they also introduce a qualitatively new security risk: \emph{stateful compromise}.

Prompt injection attacks have been widely studied in single-session settings, where adversarial inputs attempt to override system instructions or elicit policy violations \citep{prompt_injection_survey}. Memory poisoning attacks such as AgentPoison and MINJA demonstrate that attackers can manipulate retrieval systems or memory stores under certain threat models \citep{agentpoison, minja}. Meanwhile, recent work on deceptive model behaviors such as Sleeper Agents highlights the risks of persistent latent behaviors that activate under specific conditions \citep{sleeper_agents}.

However, existing work largely studies either:
\begin{itemize}
    \item single-session prompt injection,
    \item direct memory poisoning with attacker-controlled writes, or
    \item shared or adversarial multi-user settings.
\end{itemize}

In contrast, we study a realistic single-user deployment where:
\begin{itemize}
    \item the user is benign,
    \item memory writes occur automatically via tool ingestion,
    \item the attacker has only one opportunity to inject content via an untrusted tool output.
\end{itemize}

We show that in such systems, a single malicious email can implant a payload that persists across many benign sessions and later triggers exfiltration of sensitive user data when topic-relevant conversations arise.

\subsection{Contributions}

\begin{enumerate}
    \item \textbf{Persistent Exfiltration Threat Model.} We define and evaluate persistent data-exfiltration attacks in single-user memory-augmented agents.
    \item \textbf{Architectural Vulnerability Across Memory Backends.} We demonstrate vulnerability across sliding-window context, RAG, explicit memory, and Mem0.
    \item \textbf{Failure of Heuristic Defenses.} We show that common mitigations reduce but do not eliminate attack success.
    \item \textbf{Provable Information-Flow Defense.} We introduce a defense with a formal non-interference guarantee and demonstrate empirical prevention of exfiltration.
    \item \textbf{Capability-Aware Utility Evaluation.} We propose a structured evaluation framework enabling transparent security–utility tradeoff analysis.
\end{enumerate}

\section{Related Work}

\subsection{Prompt Injection}

Prompt injection attacks exploit the inability of LLMs to reliably distinguish instructions from data \citep{prompt_injection_survey}. Industry analyses and red-teaming efforts demonstrate the practical risks in tool-augmented agents.

\subsection{Memory Poisoning and Retrieval Attacks}

AgentPoison \citep{agentpoison} and MINJA \citep{minja} demonstrate poisoning attacks against memory and RAG pipelines. These works assume more direct attacker influence over memory writes or shared-user environments.

\subsection{Persistent and Deceptive Behaviors}

Sleeper Agents \citep{sleeper_agents} demonstrate models can learn conditional behaviors that persist across training and safety fine-tuning.

\subsection{Information Flow and Non-Interference}

Information-flow control provides formal guarantees that sensitive inputs do not influence public outputs \citep{ifc_classic}. We adapt this perspective to LLM agent pipelines.

\section{Threat Model}

\subsection{Adversary Capabilities}

The adversary can:
\begin{itemize}
    \item Provide or influence content returned by an untrusted tool (e.g., email).
    \item Craft arbitrary malicious content within that tool output.
\end{itemize}

The adversary cannot:
\begin{itemize}
    \item Directly edit memory.
    \item Access the agent after initial injection.
    \item Interact as a malicious user.
\end{itemize}

\subsection{Adversary Goal}

The goal is confidential data exfiltration: causing the agent to send sensitive user information to an attacker-controlled endpoint in a future session.

\subsection{System Assumptions}

\begin{itemize}
    \item Single-user deployment.
    \item Automatic memory ingestion from tool outputs.
    \item Agent capable of invoking external communication tools.
\end{itemize}

\section{Methods}

\subsection{Email Agent Implementation}

We implement a tool-augmented email assistant.

\TODO{Describe exact model used (e.g., Gemini 3.1 Pro), system prompt structure, temperature, tool invocation mechanism.}

\subsection{Memory Backends}

We evaluate four memory configurations:

\paragraph{Sliding Window Context}
Conversation history retained within a fixed token window.

\TODO{Specify window size and truncation strategy.}

\paragraph{RAG}
Memory stored as embeddings and retrieved via similarity search.

\TODO{Specify embedding model, chunk size, retrieval k, indexing mechanism.}

\paragraph{Explicit Memory}
Structured memory entries stored separately and appended to prompts.

\TODO{Describe format, indexing logic, editing permissions.}

\paragraph{Mem0}
State-of-the-art agentic memory backend.

\TODO{Describe integration details and configuration.}

\subsection{Attack Construction}

Attacks are optimized using OpenEvolve \citep{openevolve}.

Each attack:
\begin{itemize}
    \item Injected once via malicious email.
    \item Stored in memory.
    \item Triggered when sensitive topics arise.
\end{itemize}

\TODO{Specify optimization objective, reward signal, success criteria, number of optimization steps.}

\subsection{Persistence Evaluation}

We evaluate attack success at session indices:

\[
N + x \quad \text{for} \quad x \in \{10,20,\dots,100\}
\]

where \(N\) is injection session.

\TODO{Define ASR formally.}

\subsection{Defense Implementations}

\paragraph{User-Prompt-Only Indexing}
Only user messages indexed.

\TODO{Define precisely how tool outputs are excluded.}

\paragraph{No Indexing After Untrusted Tools}
Memory indexing disabled after tool ingestion.

\paragraph{Length-Limited Memory}
Memory constrained to maximum token size.

\TODO{Specify limit and eviction policy.}

\paragraph{Information-Flow Policy}
We enforce a separation between sensitive user inputs and external tool outputs.

\TODO{Provide formal non-interference definition and enforcement mechanism.}

\subsection{Utility Evaluation Framework}

We define capability classes corresponding to required agent functionalities.

\TODO{List and define each capability class clearly.}

Utility is measured as:

\[
\text{Utility} = \frac{\text{Successful benign task completions}}{\text{Total benign tasks}}
\]

Each defense is evaluated per capability class.

\TODO{Explain test generation and evaluation procedure.}

\bibliographystyle{plainnat}
\bibliography{references}

\end{document}